\subsection{Numerical Motivation and Anatomy--Function Boundary}
\label{subsec:pressure-flow-motivation}

The numerical problem in this report is to specify and audit a reduced stenosis
solver, not to infer a patient outcome. A narrowed lumen changes the local flow
field, but the same visible narrowing can have different flow consequences under
different geometry, wall, rheology, and boundary assumptions. Coronary stenosis
studies therefore distinguish anatomical descriptors from functional flow
assessment, and lesion-imaging work likewise records geometry without by itself
defining the reduced output
\cite[p.~4401]{DerimayEtAl2021CoronaryStenosisObstruction}
\cite[Abstract; Methods]{VanDenHoogenEtAl2022CoronaryLesionCTA}.

This anatomy--function distinction motivates the discipline used here: the model
state, closure choices, boundary approximation, and observation operator must be
declared before a computed quantity can be interpreted. The completed study keeps
that motivation at the level of numerical observables. It specifies an
$R_{\max}$-normalized 1D area--flow solver, tests a manufactured-solution record,
exposes the non-well-balanced geometry-rest drift, and compares velocity and physical
flow against the available resolved 3D data through a declared plane-quadrature
operator.
